% Información para el motor (arara). Compilar en LuaLaTeX o XeLaTeX sólamente.

%arara: lualatex: {
%arara: --> interaction: nonstopmode,
%arara: --> options: ['-file-line-error'],
%arara: --> shell: yes,
%arara: --> synctex: yes,
%arara: --> }
%arara: bib2gls if found("aux", "glsxtr@resource")
%arara: makeglossaries if found("aux", "@istfilename")
%arara: makeindex

\documentclass[letterpaper,twoside]{article}
%%%%%%% VARIABLES ESTÁTICAS %%%%%%%
% Aqui se definen las variables para diferentes procedimientos.
% Verificar que las unidades coincidan si se deciden cambiar.

\newcommand{\defglo}[1]{\item[{\bfseries \glsname{#1}}] \glsdesc{#1}} % Define un comando para dar un item en un entorno de descripción para una palabra definida en el glosario
% \includeonly{src/PSA-2.19.tex,src/PSA-Gloss.tex}
\newcommand{\TipoID}{TEMP} % Definirlo como: PRO, PROG, ToC, G, L, CE, CI, I, etc.
\newcommand{\Prog}{TEMP}
\newcommand{\ProgTitulo}{TEMP}

\newcommand{\ie}{\emph{i.e.}}
\newcommand{\eg}{\emph{e.g.}}

%%%%%%%%%%%%%%%% Formularios %%%%%%%%%%%%%%%%
\newcommand{\IdFormAACC}{F-OP-40}

\newcommand{\BitES}{Bitácora de entradas y salidas --- Clientes (F-AD-1)}
\newcommand{\BitULl}{Bitácora de uso de llaves (F-AC-1)}
\newcommand{\BitVig}{Bitácora de vigilancia (F-AC-2)}
\newcommand{\Fpap}{Papeleta (F-AC-3)}
\newcommand{\Oent}{Orden de entrada (F-AC-4)}
\newcommand{\Osal}{Orden de salida (F-AC-5)}
\newcommand{\LayES}{Layout de carga o descarga (F-AC-6)}
\newcommand{\BitClor}{Bitácora de cloración de agua (F-AC-7)}
\newcommand{\BitTap}{Bitácora de concentración de tapete sanitario (F-AC-8)}
\newcommand{\BitUQ}{Bitácora de uso de productos químicos (F-AC-9)}
\newcommand{\BitLAd}{Bitácora de limpieza de aduana sanitaria (F-AC-10)}
\newcommand{\BitLOf}{Bitácora de limpieza de oficinas (F-AC-11)}
\newcommand{\BitLEx}{Bitácora de limpieza de exteriores (F-AC-12)}
\newcommand{\BitBPH}{Bitácora de cumplimiento de BPH (F-HyS-1)}
\newcommand{\BitCAP}{Bitácora de consumo de agua potable (F-HyS-2)}
\newcommand{\BitCE}{Bitácora de consumo de energía eléctrica (F-HyS-3)}
\newcommand{\BitIEx}{Bitácora de inspección de extintores (F-HyS-4)}
\newcommand{\BitVIR}{Bitácora de verificación de termómetros infrarrojos (F-HyS-5)}
\newcommand{\LVCMo}{Lista de verificación de condición de montacargas (F-HyS-6)}
\newcommand{\RTAC}{Registro de temperaturas con termómetro IR --- Aseguramiento de calidad (F-HyS-7)}
\newcommand{\RTMa}{Registro de temperaturas con termómetro IR --- Mantenimiento (F-HyS-8)}
\newcommand{\RAcL}{Registro de actividad de limpieza (F-HyS-9)}
\newcommand{\RAcM}{Registro de actividad de mantenimiento (F-Hys-10)}
\newcommand{\ReEq}{Registro de revisión de equipos de enfriamiento (F-MAN-1)}
\newcommand{\RCMi}{Registro de inspección de infraestructura en prevención (F-MAN-2)}
\newcommand{\Rcap}{Registro de asistencia a capacitación (F-MAN-3)}
\newcommand{\REvDQ}{Registro de eventos de derrame de químicos (F-MAN-4)}
\newcommand{\InvVyPD}{Registro de inventario de vidrio y plástico duro (F-MAN-5)}
\newcommand{\RVerSGC}{Registro de verificación del sistema de gestión de calidad (F-MAN-6)}
\newcommand{\RRCA}{Registro de recorrido de control de alérgenos (F-MC-1)}
\newcommand{\RAC}{Registro de acción correctiva o preventiva (F-OP-1)}
\newcommand{\RACQ}{Registro de acción correctiva --- Quejas (F-OP-2)}
\newcommand{\RVSCP}{Registro de verificación del servicio de control de plagas (F-OP-3)}
\newcommand{\RevD}{Registro de revisión diaria (F-SEG-1)}
\newcommand{\RCVyPD}{Registro de evento de contaminación por VyPD (F-SEG-2)}
\newcommand{\REvCA}{Registro de evento de contaminación por alérgeno (F-SEG-3)}



%%%%%%%%%%%%%%%% CÓDIGOS PSA %%%%%%%%%%%%%%%%
% \newcommand{\PSA}{PSA-3-PRO}                 % PSA-2.3 Inspección de producto durante su almacenamiento

%%%%%%%%%%%%%%%% PSA-2.3 | Inspección de producto durante su almacenamiento %%%%%%%%%%%%%%%%
\newcommand{\TiempoAndenCong}{\qty{\leq 30}{\minute}}
\newcommand{\TiempoAndenRefri}{\qty{\leq 30}{\minute}}
\newcommand{\TiempoAveriaRefri}{4}              % Tiempo máximo de permanencia de alimento refrigerado en cámara averiada (en horas) 
\newcommand{\TiempoAveriaConge}{10}             % Tiempo máximo de permanencia de alimento congelado en cámara averiada (en horas)
\newcommand{\EspacioMinimoParedProducto}{30}    % Espacio entre tarima y pared (en centimetros)
\newcommand{\EspacioMinimoTechoProducto}{30}    % Espacio entre tarima y techo (en centimetros)
\newcommand{\VecesTempManual}{dos}              % Veces que se verifica la temperatura manualmente



%%%%%%%%%%%%%%%% PSA-2.3 | Política de aprobación de proveedores %%%%%%%%%%%%%%%%
\newcommand{\NotifAuditoriaAProveedor}{2 semanas}           % Plazo máximo para notificar a proveedor sobre auditoría.
\newcommand{\PlazoResultadoAProveedor}{2 semanas}           % Plazo máximo para notificar a proveedor sobre resultados de auditoría.
\newcommand{\PlazoPlanDeAccionAuditProveedor}{2 semanas}    % Plazo maximo para recibir plan de acción de parte del proveedor tras auditoría

%%%%%%%%%%%%%%%% PSA-2.15 | Almacenamiento de registros %%%%%%%%%%%%%%%%
\newcommand{\VigenciaAlmacRegistros}{2 años}
\newcommand{\VigenciaAlmacRegistrosElec}{3 años}

%%%%%%%%%%%%%%%% GLOBAL | Resetear glosarios en cada sección %%%%%%%%%%%%%%%%
\usepackage{etoolbox}
\preto\section{\glsresetall}

\newcommand{\AC}{aseguramiento de calidad}
\newcommand{\G}{gerencia}
\newcommand{\AD}{alta dirección}
\newcommand{\OP}{operaciones}
\newcommand{\Emb}{embarques}
\newcommand{\MC}{mesa de control}	% Definición de variables estáticas
								% i.e. Tiempo máximo de cámras de refrigeración sin energía

%%%%%%%%%%%%%%%%%%%%%%%%%%% ESPECIALES %%%%%%%%%%%%%%%%%%%%%%%%%%%%%%%%%%
% \RequirePackage{expl3}    % \usepackage cannot be used before \documentclass
% \ExplSyntaxOn             % Switch on expl3 syntax
% \pdf_version_gset:n{1.7}  % Use provided expl3 function
% \ExplSyntaxOff            % Switch off expl3 syntax

%%%%%%%%%%%%%%%%%%%%%%%%%%% Tipografía %%%%%%%%%%%%%%%%%%%%%%%%%%%%%%%%%%

\usepackage[sfdefault,t,lf]{carlito}    % Alternativa OpenSource a Calibri, definir sans serif como predeterminado, numerales tabulares
% \usepackage[default,lnum]{opensans}
% \usepackage{CormorantGaramond}

% \usepackage{lucidabr}

% \usepackage{fontspec}

\fontsize{8.8pt}{12pt}\selectfont
\renewcommand{\baselinestretch}{1.2}

% \linespread{1.3}
% \DeclareUnicodeCharacter{2212}{-}

\usepackage{polyglossia}
	\setdefaultlanguage[spanishoperators=all]{spanish}  % Definir la localización del documento como es-MX (para que cuadro sea tabla, etc.)
	% \setdefaultlanguage{mexican}
	\renewcommand{\listtablename}{Índice de tablas}

\usepackage{changelog}  % Control de versiones

\usepackage{pdflscape}
\usepackage{pifont}
\newcommand{\cmark}{\ding{51}}%
\newcommand{\xmark}{\ding{55}}%

\usepackage[
	record=hybrid,% provides group and location fields (and other stuff)
	nopostdot,
	xindy,
	acronym,
	stylemods=tree]{glossaries-extra}
\setabbreviationstyle{long-short-sc}
\GlsXtrLoadResources[
	src={../glosario.bib},	% bib files
	selection={all},			% select all entries
	sort={es-MX}
]
\makeglossaries

\renewcommand{\glossarypreamble}{\fontsize{10pt}{12pt}\selectfont}
	
\usepackage{imakeidx}       % Para hacer indices
	\makeindex[columns=2]   % Estilo del indice alfabético

\newcounter{note}[section]  % Nuevo contador para notas
\renewcommand{\thenote}{\thesection.\Alph{note}}    % Definir estilo de numeración de las notas
\usepackage[framemethod=TikZ]{mdframed}             % Motor gráfico para notas (TikZ)
\newenvironment{note}[1][]{%                        % Nuevo entorno de Nota (note) 
    \refstepcounter{note}
    \begin{mdframed}[%                                  
        frametitle={Nota \thenote\ #1},                 % Titulo de las notas
        skipabove=\baselineskip{} plus 2pt minus 1pt,     % Espaciado superior
        skipbelow=\baselineskip{} plus 2pt minus 1pt,     % Espaciado inferior
        linewidth=0.5pt,                                % Tamaño de linea del recuadro
        frametitlerule=true,                            % incluir titulo
        frametitlebackgroundcolor=gray!30               % Color del recuadro
    ]%
}{%
    \end{mdframed}
}

\newcounter{contact}[section]                           % Nuevo contador para contactos
\renewcommand{\thecontact}{\thesection.\Alph{contact}}  % Definir que la numeración sea Sección.Alfabético

\newenvironment{contact}[1][]{%
    \refstepcounter{contact}
    \begin{mdframed}[%
        frametitle={Contacto \thecontact\ #1},          % Titulo del contacto
        skipabove=\baselineskip{}plus 2pt minus 1pt,	% Espaciado superior
        skipbelow=\baselineskip{}plus 2pt minus 1pt,    % Espaciado inferior
        linewidth=0.5pt,                                % Tamaño de linea del recuadro
        frametitlerule=true,                            % poner titulo
        frametitlebackgroundcolor=yellow!30             % Color del recuadro
    ]%
}{%
    \end{mdframed}
}

\usepackage[multiple]{footmisc}     % Cambios en la presentación de pies de página
\usepackage{microtype}              % Ajustes de kerning y otros 
\usepackage{graphicx}               % Mejora al motor de imagenes
\usepackage{hyperref}               % Referencias cruzadas en PDF
\usepackage{cleveref}               % Referencias cruzadas en texto
%%%%%% REDEFINICIONES DE REFERENCIAS CRUZADAS %%%%%%%%%
    \crefname{note}{nota}{notas}    
	\Crefname{note}{Nota}{Notas}
	\crefname{contact}{contacto}{contactos}
	\Crefname{contact}{Contacto}{Contactos}
	\crefname{paragraph}{parrafo}{parrafos}
	\Crefname{paragraph}{Parrafo}{Parrafos}
%%%%%%%%%%%%%%%%%%%%%%%%%%%%%%%%%%%%%%%%%%%%%%%%%%%%%%%
\newcommand{\email}[1]{\href{mailto:#1}{#1}}    % Comando para hacer que se puedan clickear emails
\newcommand{\cels}[1]{\qty{#1}{\celsius}}       % Escrir celsius rapido

\usepackage{lastpage}   % Definir variable para el numero total de páginas
\usepackage{siunitx}    % Unidades internacionales y alineamiento de decimales en tablas
	\sisetup{detect-all}    % Detectar automaticamente que la tipografía es sans serif

\usepackage{color}  % definir colores nuevos

\usepackage{tabularray} % Creación de tablas
	\UseTblrLibrary{booktabs,siunitx}   % librerias parea tablas correctas, unidades internacionales
	\definecolor{Gallery}{rgb}{0.929,0.929,0.929}   % definición del color "Gallery"
	\DefTblrTemplate{contfoot-text}{normal}{\textit{Continúa en la siguiente página}}   % Continuación de tabla
	\SetTblrTemplate{contfoot-text}{normal} % Continuación de tabla normal (footer)
	\DefTblrTemplate{conthead-text}{normal}{\textit{(Continuación)}}    % Continuación de tabla
	\SetTblrTemplate{conthead-text}{normal} % Continuación de tabla (header)

    % \usepackage[margin=1cm,bmargin=2cm,tmargin=5cm,headheight=110pt]{geometry}   % Dispoción espacial del documento
% \newgeometry{includeheadfoot, heightrounded, left=1cm, right=1cm, top=1cm, bottom=4cm, headheight=6em, footskip=1.5em}    
%     \savegeometry{L1}

% \newgeometry{includeheadfoot, heightrounded, left=1cm, right=1cm, top=1cm, bottom=4.5cm, headheight=7em, footskip=1.5em, showcrop, showframe}
%     \savegeometry{L2}

\usepackage[%
includeheadfoot,
heightrounded,
left=1cm,
right=1cm,
top=1cm,
bottom=1cm,
headheight=110pt,
footskip=1.5em]{geometry}    

\usepackage[export]{adjustbox}

\usepackage{enumitem}
	\setlist[1]{labelindent=\parindent} % Indentación como parrafo para listas
	\setlist[description]{style=nextline}
	\setlist[itemize]{leftmargin=*} 	% Eliminar margen en izq
	\setlist[itemize,1]{label=---}  	% Viñeta como em-dash
	\setlist[itemize,2]{label=--}
	\setlist{noitemsep} 				% Que no se separen los item
	% [label={\theparagraph--\arabic*}] % Numerarar parrafo-numero

	\newlist{ListaSec}{enumerate}{3}
	\setlist[ListaSec]{label* = \arabic*., name* = \arabic{section}.\arabic{subsection}}


%%%%%%%%%%%%%%%% DEFINICIÓN DE CABECERA %%%%%%%%%%%%%%%%
\usepackage{fancyhdr}   % Usar fancy headers y footers
	\renewcommand{\headrulewidth}{0pt}  % Grosor de regla del header
	\fancyhead[CE,CO,LE,LO,RE,RO]{} % Resetear headers    
    \fancyfoot[CE,CO,LE,LO,RE,RO]{} % Resetear footers

%%%%%%%%%%%%%%%%%%%%%% VARIABLES DE CADA DOCUMENTO DENTRO DE UN PROGRAMA %%%%%%%%%%%%%%%%%%%%%
	\newcommand{\Codigo}{}                  % Codigo del documento
	\newcommand{\FechaPub}{}                % Fecha de publicación
	\newcommand{\Titulo}{}                  % Titulo del documento
	\newcommand{\MayorVer}{}                % Versión mayor
	\newcommand{\MenorVer}{}                % Verisón menor
	\newcommand{\Edit}{\MayorVer.\MenorVer} % Definición de la edición

    \newcommand{\Elaboro}{}                 % Nombre de quien lo elaboró
    \newcommand{\Reviso}{}                  % Nombre de quien revisó el documento
    \newcommand{\Autorizo}{}                % Nombre de quien aprobó el documento

%%%%%%%%%%%%%%%%%% Formato por defecto para informaicón documentada general %%%%%%%%%%%%%%%%%%
\fancypagestyle{formato-PS}{%
	\fancyhead{}
	\fancyfoot{}
	\fancyhead[RE]{%
		\begin{tblr}{%	
			colspec = {X[c,m]X[c,m]},    % Especificación de columnas: Anchura estática, centrado en medio
			hlines,
			vlines,
			width = 0.35\linewidth,}%
			{\textbf{Código:}\\ \Codigo} 				 & {\textbf{Edición:}\\ \MayorVer.\MenorVer} 			\\ 	% Logotipo (2 columnas), red de frios (3 columnas), código del documento
			{\textbf{Fecha de publicación:}\\ \FechaPub} & {Página~\thepage~de~\pageref{LastPage}} 		% Titulo (4 columnas), Fecha, edición
		\end{tblr}
	}%

	\fancyhead[LO]{%
		\begin{tblr}{%	
			colspec = {X[c,m]X[c,m]},   % Especificación de columnas: Anchura estática, centrado en medio
			hlines,
			vlines,
			width = 0.35\linewidth,}%
			{\textbf{Código:}\\ \Codigo} 					& {\textbf{Edición:}\\ \MayorVer.\MenorVer} 			\\	% Logotipo (2 columnas), red de frios (3 columnas), código del documento
			{\textbf{Fecha de publicación:}\\\FechaPub}	& {Página~\thepage~de~\pageref{LastPage}} 		% Titulo (4 columnas), Fecha, edición
		\end{tblr}
	}%

	\fancyfoot[C]{Página~\thepage~de~\pageref{LastPage}}   % Para que muestre Página # de #
	\fancyfoot[L]{\textbf{CONFIDENCIAL}}                    % CONFIDENCIAL
	\fancyfoot[R]{\Codigo~v.\Edit}                          % Estructura del código CODIGO v.EDICION
}%

\fancypagestyle{formato-PI}{%
	\fancyhead{}
	\fancyfoot{}
    \fancyhead[C]{%
	    \begin{tblr}{%	
            colspec = {X[c,m]X[c,m]X[c,m]X[c,m]X[c,m]X[c,m]},    % Especificación de columnas: Anchura estática, centrado en medio
            hlines,
            vlines,
			width = \linewidth,}%
			\SetCell[c=2]{} \includegraphics[height=3em, valign=c]{./RDF_Logo.png}	&	& \SetCell[c=3]{} \large{Red de fríos S.A. de C.V.}	&	&												& {\textbf{Código:}\\ \Codigo}	\\ % Logotipo (2 columnas), red de frios (3 columnas), código del documento
			\SetCell[c=4]{} {\textbf{Titulo:}\\ \Titulo} 							&	&													&	& {\textbf{Fecha de publicación:}\\ \FechaPub}	& {\textbf{Edición:}\\ \Edit}	\\ % Titulo (4 columnas), Fecha, edición
            \end{tblr}
	}
    \fancyfoot[C]{Página~\thepage~de~\pageref{LastPage}}   % Para que muestre Página # de #
    \fancyfoot[L]{\textbf{CONFIDENCIAL}}                    % CONFIDENCIAL
    \fancyfoot[R]{\Codigo~v.\Edit}                          % Estructura del código CODIGO v.EDICION
}%

%%%%%%%%%%%%%%%%%%%%%%%%%%%%%%%%%%%%%%%%%%%%%%%%%%%%%%%%%%%%%%%%%%%%%%%%%%%%%%%%%%%%%%%%%%%%%%%%%%%%%%%%%%%%%%%%%%
			
\setcounter{tocdepth}{4}            % Qué tanto mostrará la tabla de contenidos
\setcounter{secnumdepth}{6}         % Que tantas secciones numerará LaTeX (Se ajustó hasta subparrafo)

 % Define un comando para dar un item en un entorno de descripción para una palabra definida en el glosario

% \includeonly{src/BPM-1.1.tex}
\renewcommand{\Prog}{ALE}
\renewcommand{\ProgTitulo}{Programa de control de alérgenos}

\begin{document}
\pagestyle{formato-PS}
\thispagestyle{formato-PI}
\renewcommand{\MayorVer}{1}
\renewcommand{\MenorVer}{0}
\renewcommand{\TipoID}{ToC}
\renewcommand{\Codigo}{\Prog--\TipoID}
\renewcommand{\FechaPub}{2023--01}
%\renewcommand{\Edit}{01}
\renewcommand{\Titulo}{\ProgTitulo --- Índice}

\tableofcontents
\thispagestyle{formato-PI}
\renewcommand{\MayorVer}{2}
\renewcommand{\MenorVer}{0}
\renewcommand{\Titulo}{Programa de control de alérgenos}
\renewcommand{\TipoID}{ESP}
\renewcommand{\FechaPub}{2023--01}

\section{\Titulo}\index{Programa!control de alérgenos, de}\index{Especificación!capacitaciones para el personal del almacén}
\renewcommand{\Codigo}{\Prog--\thesection--\TipoID}

\subsection{Objetivo}
Establecer un procedimiento el correcto manejo de los productos y materiales alérgenos y evitar la contaminación cruzada con otros productos almacenados.

\subsection{Alcance}
Alimentos que contengan componente(s) alérgeno(s).

\subsection{Terminología y definiciones}
\begin{description}
	\defglo{peligro-relacionado-con-la-inocuidad-de-los-alimentos};
	\defglo{contacto-cruzado-alergenos};
	\defglo{alergeno}.
\end{description}

\subsection{Documentos y/o normas relacionadas}
\begin{itemize}
	\item Programa de Higiene y Sanidad.
	\item 27 CFR 7.82 --- Voluntary disclosure of major food alergens.
	\item 21 CFR 117.3 --- Definitions.
	\item Código SQF 2.8.1.
\end{itemize}

\subsection{Control de alérgenos en planta}
\begin{itemize}
	\item Alguno de los productos almacenados de \gls{RDF} son productos preparados formulados y contienen uno a varios productos alérgenos, cabe mencionar que estos productos se encuentran envasados (Empaque original) y emplayados.
	\item Estos productos que contienen ingredientes alérgenos son manejados de la siguiente manera:
	\item Áreas específicas y señaladas con la leyenda \emph{ALERGENO} para almacenamiento de productos que contienen alérgenos.
	\item Cada producto es almacenado por separado y en ningún momento son mezclados.
	\item Independientemente del programa interno de alérgenos algunos de estos productos se encuentran identificados de origen con la leyenda \emph{ALERGENO}
	\item Cada producto DEBE de estibarse con productos de su mismo tipo, tarimas del mismo producto.
	\item Registrar todos los movimientos de estos productos (bitácora de entrada y salida).
	\item Identificar en cada uno de los puntos del proceso con etiqueta de alérgenos, recepción, alimentación y embarque.
\end{itemize}

\subsubsection{Contacto cruzado de alérgenos}
La clave para manejar alérgenos en el procedimiento y manejo, es evitar el \gls{contacto-cruzado-alergenos} Si el mismo ingrediente alergénico fuese utilizado en todas las fórmulas de producto, entonces no habría riesgo de contacto-cruzado usualmente. Se deben utilizar políticas y procedimientos que deben incluir la documentación apropiada para apoyar estas actividades. Elementos clave que deben ser considerados al desarrollar e implementar estas políticas y procedimientos que incluyen.

\begin{itemize}
	\item Limpieza durante cambio de alérgenos
	\item Inspecciones de almacén.
	\item Colocación de productos alérgenos sobre productos alérgenos del mismo tipo.
	\item Barreras físicas laterales como empaque o emplaye.
	\item Colocar o ubicar en el rack productos no alérgenos sobre productos alérgenos.
	\item Códigos de color u otra designación y segregación de tarimas.
	\item Áreas exclusivas.
	\item Todos los productos de este tipo deben ser inspeccionadas por daños (incluyendo fugas, cajas dañadas, etc.). Cualquier producto que se encuentre con fugas DEBEN ser devueltos inmediatamente y documentar en él la bitácora de “Registro de Producto Dañado”. Los esfuerzos de limpieza deben garantizar que no existe contaminación cruzada con otros productos.
	\item La limpieza con agua se recomienda para eliminar cualquier residuo pastoso o pegajoso que contenga material alérgeno. Cuando se utilice un sistema de limpieza en las áreas de operación debe ser examinado para evidenciar áreas que no puedan ser adecuadamente limpiadas y que se puedan atrapar residuos alergénicos.
	\item Tanto la limpieza con agua como la limpieza en seco deben ser validadas periódicamente a través de una revisión visual para verificar que las superficies estén limpias.
	\item Después de una validación inicial del procedimiento de limpieza, se puede usar la examinación visual y diaria de superficies en contacto con el producto para verificar que la limpieza de alérgenos se ha llevado a cabo.
	\item Esta examinación visual debe documentarse en la bitácora destinada para instalaciones en las cuales se utilizó el material alérgeno.
\end{itemize}

\subsection{Responsables de la actividad}
\begin{itemize}
	\item \textbf{Ejecutado} por personal de calidad.
\end{itemize}

\subsection{Acciones correctivas}
\begin{itemize}
	\item Cuando se presente una no conformidad en la realización de los procedimientos, detectada por el supervisor o encargado del área se deberá de repetir este mismo
	\item En caso de que la desviación sea mayor, se deberá de registrar la acción correctiva en el formato “Registro de reporte y acciones correctivas”
	\item Dar aviso al supervisor inmediato.
\end{itemize}

\subsection{Frecuencia}
\begin{itemize}
	\item Cada producto que contiene alérgeno.
\end{itemize}

\begin{changelog}[title=Registro de cambios,simple, sectioncmd=\subsection*,label=changelog-\thesection-\MayorVer.\MenorVer]
	\begin{version}[v=\MayorVer.\MenorVer, date=2023--01, author=Pablo E. Alanis]
		\item Cambio de formato;
		\item Cambios en la serialización de versiones;
		\item Correcciones ortográficas y de estilo.
	\end{version}

	\begin{version}[v=1.7, date=2022--05, author=Alonso M.]
		\item cambio de fecha;
	\end{version}

	\shortversion{v=1.6, date=2021--05, changes=No hubo cambios}
\end{changelog}

\thispagestyle{formato-PI}
\renewcommand{\MayorVer}{2}
\renewcommand{\MenorVer}{1}
\renewcommand{\Codigo}{PSA-1-PROG} %TODO
\renewcommand{\FechaPub}{2023--01}
%\renewcommand{\Edit}{2.1}
\renewcommand{\Titulo}{Manejo de proteína cruda fresca}

\section{\Titulo}
\label{PRO-ManejoDeProtFresca}
\index{Información documentada!tipo!procedimiento!Manejo de proteína cruda fresca}

% \section{Manejo de proteína cruda fresca}


\subsection{Objetivo}

Asegurar que los productos de proteína fresca cruda se almacenan de manera que no sean causa de contaminación cruzada.

\subsection{Alcance}

Todo producto que contenga proteína cruda fresca.

\subsection{Términos y definiciones}

\begin{itemize}
	\item \textbf{Proteína cruda:} Se le llama así al análisis químico que se realiza en los alimentos, basados en el nitrógeno, para obtener la cantidad de proteínas que este contiene.
\end{itemize}

\subsection{Documentos y/o normas relacionadas}

\begin{itemize}
	\item Programa de Sanidad
	\item Programa de Alérgenos.
\end{itemize}

\subsection{Procedimiento}

\subsubsection{Precauciones de seguridad}

\begin{itemize}
	\item Usar uniforme completo.
\end{itemize}

\subsubsection{Instrucciones}

\paragraph{Almacenaje de proteína cruda}

\begin{itemize}
	\item Todas las proteínas crudas frescas deben de ser almacenadas en congelación, estas deben de almacenarse de manera que no sean causa de contaminación cruzada con otros productos o materiales de empaque.
	\item No deben almacenarse sobre productos listos para ser consumidos (RTE).
	\item No deben de almacenarse sobre materiales de empaque u otros productos como pudieran ser ingredientes.
	\item Los productos que contengan proteína cruda fresca deben de almacenarse de forma tal que en caso de que ocurriera un derrame no contaminen otros productos, esto quiere decir que solo se deben de almacenar:
	\begin{itemize}
		\item El producto con proteína cruda fresca se debe de almacenar en la parte más baja del rack y sobre él, los productos que no contengan proteína cruda.
		\item Por ningún motivo debe de almacenarse producto de empaque o producto RTE por debajo de productos que contengan proteína cruda fresca.
		\item El producto con proteína cruda fresca se debe de almacenar sobre producto con proteína cruda.
		\item Se debe de almacenar productos con proteína cruda fresca sobre producto del mismo tipo de proteína cruda.
		\item Todo el producto debe estar paletizado para generar barrera física con productos almacenados de manera lateral.
		\item En caso de derrames revisar procedimiento de limpieza para proteínas crudas frescas.
	\end{itemize}
\end{itemize}

\paragraph{Embarque de proteína cruda}

\begin{itemize}
	\item Toda tarima con producto que contenga proteína cruda fresca debe ser paletizado como barrera física y evitar contaminación cruzada con productos almacenados de manera lateral.
	\item En el caso de encontrarse varias tarimas con proteína cruda fresca, estas deben de colocarse juntas al momento del embarque para minimizar la cercanía con otros productos.
	\item Cuando se tienen tarimas compuestas (varios productos), las proteínas crudas frescas deben de ser acomodadas en la parte baja de la tarima y estos productos deben de ser paletizado como barrera física antes de colocar el resto de los productos.
	\item Por ningún motivo debe de almacenarse producto de empaque o producto RTE por debajo de productos que contengan proteína cruda fresca.
	\item El producto con proteína cruda fresca se debe de almacenar sobre producto con proteína cruda fresca.
\end{itemize}


\paragraph{Alérgenos}

\begin{itemize}
	\item Algunos de las proteínas crudas frescas también son consideradas como alérgenos como ejemplos se pueden mencionar el huevo, mariscos, los cuales se deben de manejar en base a lo indicado a este procedimiento y al procedimiento de manejo de productos Alérgenos.
\end{itemize}

\subsection{Responsables de la actividad}

\begin{itemize}
	\item \textbf{Ejecutado} por personal de operaciones
	\item \textbf{Monitoreado} por personal de calidad
	\item \textbf{Verificado} por personal de gerencia.
\end{itemize}

\subsection{Acciones preventivas}

\begin{itemize}
	\item Se llevara a cabo una inspección. Registrando el indicador y medida correspondiente
	\item Si la desviación se repite frecuentemente se dará curso de capacitación al personal de almacén y limpieza para que realicen eficientemente su trabajo.
	\item Si después de haber capacitado al personal de almacén se siguen presentando desviaciones por causas injustificadas, se tomaran acciones más enérgicas con el personal por incumplimiento con sus deberes.
\end{itemize}

\subsection{Acciones correctivas}

\begin{itemize}
	\item En caso de no cumplimiento la tarea se deberá volver a realizar como se indica en el procedimiento.
	\item En caso de no conformidad reportar en formato de acciones correctivas F-OP-40.xls|F-OP-40
\end{itemize}

\subsection{Frecuencia}

\begin{itemize}
	\item Cada recepción de producto.
	\item Cada entrega de producto.
\end{itemize}

\subsection{Historial de modificaciones}

\begin{itemize}
	\item \textbf{Cuarta edición:} cambio de fecha de 23 de octubre 2017 a 28 de enero 2019, se realizó cambio de código de DG-P-MPC a PR-016. Y no se realizaron cambios de la revisión 003 a la 004.
	\item \textbf{Quinta edición:} febrero 2020 se hizo cambio de formato y cambio de código de PR-016 a PRO-OP-012. Y no se realizaron cambios de la revisión 04 a la 05.
	\item \textbf{Sexta edición:} febrero 2021 no se realizaron cambios de la revisión 05 a la 06.
	\item \textbf{Séptima edición:} febrero 2022 no se realizaron cambios de la revisión 06 a la 07.
\end{itemize}

\subsection{Listado de distribución}


\thispagestyle{formato-PI}
\renewcommand{\MenorVer}{0}
\renewcommand{\MayorVer}{2}
%\renewcommand{\Edit}{\MayorVer.\MenorVer}
\renewcommand{\Codigo}{HYS-13-IT}
\renewcommand{\FechaPub}{2023--01}
\renewcommand{\Titulo}{Limpieza de utensilios de limpieza de alérgenos}

\section{\Titulo}\index{Limpieza de utensilios!limpieza de alérgenos, de}
% \section{Limpieza de accesorios de alérgenos}

\subsection{Objetivo}
Establecer un procedimiento que aseguren un mantenimiento y limpieza apropiada en caso de presentarse un derrame de producto Alérgeno.

\subsection{Alcance}
Toda superficie afectada por un derrame de alimento con componente(s) alérgenico(s).

\subsection{Terminología y definiciones}
\begin{itemize}
	\item \textbf{Limpieza:} Proceso de remover los residuos de Orgánicos y suciedad que pueden ser una fuente de contaminación. Los métodos de limpieza adecuados y los materiales dependerán de la naturaleza del alimento.
	\item \textbf{Sanitización:} Proceso realizado después de efectuar la limpieza profunda de equipos, maquinaria e instalaciones la cual nos permite el reducir el número de microorganismos
\end{itemize}
\subsection{Procedimiento}
\subsubsection{Materiales}
\begin{itemize}
	\item Preparación de \emph{Liquid K} a la concentración indicada por el proveedor, \qtyrange{0.13}{0.26}{\percent} o bien \qtyrange{1.3}{2.6}{\milli\liter\per\liter};
	\item Cubeta de \qty{20}{\litre};
	\item Cepillo para piso color Blanco;
	\item Cepillos manuales color Blanco;
	\item Recogedor de plástico;
	\item Jalador color Blanco;
	\item Toallas desechables.
\end{itemize}

\subsubsection{Instrucciones}
\paragraph{Derrame de alérgenos en el piso de almacén, cámaras frías y andenes}

\begin{enumerate}
	\item Verificar la no afectación de ningún otro producto
	\item Recolectar el producto alérgeno de las áreas afectadas con ayuda de toallas desechables (elimine toallas después de su uso)
	\item Prepare detergente a la concentración indicada por el proveedor.
	\item Tallar el piso con agua y jabón y el cepillo asignado para estos casos
	\item Enjuagar el piso
	\item Eliminar la humedad con el jalador asignado
	\item Dar aviso al supervisor de Almacén y Calidad para realizar revisión visual
	\item Registro de lavado por derrame de material alérgeno.
	\item Si se afectó algún otro producto o empaque se procede a eliminar el producto contaminante por medio de limpieza en seco (toallas desechables).
	\item Separar el producto y etiquetar como producto detenido y/o retenido
	\item Valorar el grado de afectación del producto en coordinación con el Coordinador de calidad del cliente en cuestión.
	\item Se precederá según indicaciones sobre la disposición del producto.
	\item Registro de lavado por derrame de material alérgeno.
\end{enumerate}

\subsection{Responsables de la actividad}
\begin{itemize}
	\item \textbf{Ejecutado} por personal de limpieza;
	\item \textbf{Monitoreado} por personal de mantenimiento;
	\item \textbf{Verificado} por personal de Calidad.
\end{itemize}

\subsection{Acciones preventivas}
\begin{itemize}
	\item Se llevará a cabo una inspección diaria. Registrando el indicador y medida correspondiente
	\item Si la desviación se repite frecuentemente se dará curso de capacitación al personal de almacén y limpieza para que realicen eficientemente su trabajo.
	\item Si después de haber capacitado al personal de almacén y limpieza se siguen presentando desviaciones por causas injustificadas, se tomaran acciones más enérgicas con el personal por incumplimiento con sus deberes.
\end{itemize}

\subsection{Acciones correctivas}
\begin{enumerate}
	\item La tarea es liberada una vez que el área se encuentran libre de polvo, suciedad, residuos de materiales extraños, residuos de jabón, y seca. En caso contrario la tarea se deberá volver a realizar como se indica en el procedimiento.
	\item En caso de no conformidad reportar en \RAC.
\end{enumerate}

\subsection{Frecuencia}
Cuando se presente la incidencia.

\begin{changelog}[simple, sectioncmd=\subsection*,label=changelog-3.13]
	\begin{version}[v=2.0, date=2023--01, author=Pablo E. Alanis]
		\item Cambio de formato;
		\item Cambios en la serialización de versiones;
		\item Correcciones ortográficas y de estilo.
	\end{version}

	\begin{version}[v=1.7, date=2022--05, author=Alonso M.]
		\item cambio de fecha;
	\end{version}	

	\shortversion{v=1.6, date=2021--05, changes=No hubo cambios}
\end{changelog}
\renewcommand{\MayorVer}{1}
\renewcommand{\MenorVer}{0}
\renewcommand{\TipoID}{G}
\renewcommand{\Codigo}{\Prog-\TipoID}
\renewcommand{\FechaPub}{2023--01}
%\renewcommand{\Edit}{01}
\renewcommand{\Titulo}{\ProgTitulo --- Glosario}

\printunsrtglossary[type=acronym,nonumberlist]
\printunsrtglossary[type=\glsdefaulttype,nonumberlist]
% \printunsrtglossary[type=main]    
% \printglossary[type=acronym]
% \printglossary[type=main]
% \printunsrtglossary[type=acronym]
% \printunsrtglossary[type=main]
\renewcommand{\MayorVer}{1}
\renewcommand{\MenorVer}{0}
\renewcommand{\TipoID}{I}
\renewcommand{\Codigo}{\Prog-\TipoID}
\renewcommand{\FechaPub}{2023--01}
%\renewcommand{\Edit}{01}
\renewcommand{\Titulo}{\ProgTitulo --- Indice alfabético}
\printindex
\end{document}